\subsection{星代数}

定义 3. --- 星代数（C*-代数）是这样的对合 Banach 代数 A：\(\|x\|^{2}=\|x^{*}x\|\) 对一切 \(x\in A\) 成立。

若 A 与 B 是星代数，从 A 到 B 的态射，或称星代数态射，是指从 A 到 B 的对合代数态射。从 A 到 B 的同构是指从 A 到 B 的对合代数同构。

某些作者称之为“C*-代数”。

设 A 是一个星代数。A 的星子代数是指 A 的一个闭对合子代数。

例. --- 1) 复 Hilbert 空间的连续自同态组成的对合 Banach 代数 (I, p.~100, 例 3) 是一个星代数 (EVT, V, p.~39, 命题 2)。

\begin{enumerate}
\def\labelenumi{\arabic{enumi})}
\setcounter{enumi}{1}
\tightlist
\item
  设 X 是局部紧拓扑空间。X 上连续有界复值函数组成的对合 Banach 代数 \(\mathcal{C}_{b}(X)\) (I, p.~100, 例 1) 是一个星代数。事实上，对每个函数 \(f \in \mathcal{C}_b(\mathrm{X})\)，有 \(f^*f = |f|^2\)，从而 \(\| f^*f \| = |||f|^2 \| = \| f\|^2\)。
\end{enumerate}

设 \(\mathrm{A} = \mathcal{C}_0(\mathrm{X})\) 是在无穷远处趋于 0 的函数组成的对合 Banach 子代数。它是 \(\mathcal{C}_b(\mathrm{X})\) 的一个星子代数。对每个函数 \(f \in \mathrm{A}\)，有 \(\| f \| = \varrho(f)\)，因为 \(\mathrm{Sp}_{\mathrm{A}}'(f) = f(\mathrm{X}) \cup \{0\}\)。

设 X 与 Y 是局部紧拓扑空间。对每个从 X 到 Y 中的真偏映射 \(\varphi\) (I, p.~33, 定义 1)，代数态射 \(\varphi^{*}\)（从 \(\mathcal{C}_{0}(\mathrm{Y})\) 到 \(\mathcal{C}_{0}(\mathrm{X})\)，I, p.~34, 命题 3）是对合代数态射。反之，每个星代数态射 \(\pi\colon\mathcal{C}_{0}(\mathrm{Y})\to\mathcal{C}_{0}(\mathrm{X})\) 都具有这种形式（同前）。

\begin{enumerate}
\def\labelenumi{\arabic{enumi})}
\setcounter{enumi}{2}
\item
  设 X 是紧拓扑空间，\(x_{0} \in X\) 是 X 的一个固定元素。\(\mathcal{C}'(X)\) 是 \(\mathcal{C}(X)\) 的一个由连续函数 \(f: X \to C\)（其中 \(f(x_{0}) = 0\)）组成的子代数，它是一个星代数。
\item
  设 X 是分离拓扑空间，\(\mu\) 是 X 上的一个正测度。对合 Banach 代数 \(\mathrm{L}^{\infty}(\mathrm{X},\mu)\) 是一个交换的幺星代数。
\item
  设 \((\mathbf{A}_{i})\) 是一族星代数。\(A_{i}\) 的乘积对合 Banach 代数 A (I, p.~101, 例 5) 是一个星代数，称为 \(A_{i}\) 的乘积星代数。
\item
  设 A 是一个星代数。若 B 是 A 的一个闭对合子代数，则 B 是一个星代数。（V，将出版）中将看到，每个星代数都同构于一个 Hilbert 空间的自同态组成的星代数（例 1）的一个闭对合子代数。
\item
  设 A 是一个星代数。若 M ⊂ A 是任意子集，则 A 的由 M 生成的闭对合子代数是一个星代数，称为 A 的由 M 生成的星子代数。此外若 A 是幺的，则由 M 生成的闭幺对合子代数是一个幺星代数，称为 A 的由 M 生成的幺星子代数。
\item
  一般而言，对合 Banach 代数 \(\mathcal{M}^{1}(G)\) (I, p.~100, 例 4) 不是星代数 (I, p.~181, 习题 8)。
\end{enumerate}

引理 6. --- 设 A 是一个 Banach 代数，装备了一个满足

\[
\| x \| ^ {2} \leqslant \| x ^ {*} x \|\tag{1}
\]

（对一切 \(x \in \mathbf{A}\)）的对合。那么 A 是一个星代数。

设 \(x \in \mathbf{A}\)。于是有 \(\|x\|^2 \leqslant \|x^*\| \cdot \|x\|\)，从而 \(\|x\| \leqslant \|x^*\|\)。交换 \(x\) 与 \(x^*\) 的角色，我们得到 \(\|x\| = \|x^*\|\)。于是 \(\|x^*x\| \leqslant \|x^*\|\|x\| \leqslant \|x\|^2\)，而假设蕴含等式 \(\|x\|^2 = \|x^*x\|\)。

引理 7. --- 设 A 是一个星代数。

\begin{enumerate}
\def\labelenumi{\alph{enumi})}
\tightlist
\item
  A 的正则表示 \(\gamma\) (I, p.~16, 定义 1) 是等距的，也就是说，
\end{enumerate}

\[
\| x \| = \sup _ {\| y \| \leqslant 1} \| x y \|,
\]

对一切 \(x\in \mathbf{A}\) 成立。

\begin{enumerate}
\def\labelenumi{\alph{enumi})}
\setcounter{enumi}{1}
\tightlist
\item
  对每个 \(x \in A_{h}\)，有
\end{enumerate}

\[
\varrho (x) = \| x \|.\tag{2}
\]

设 \(x \in \mathrm{A}\)。我们有 \(\sup_{\| y \| \leqslant 1} \| xy \| \leqslant \| x \|\)。为证 \(\| x \| \leqslant \sup_{\| y \| \leqslant 1} \| xy \|\)，可设 \(\| x \| = 1\)。于是元素 \(y = x^*\) 满足 \(\| y \| = \| x^* \| = 1\)，且 \(\| xy \| = \| x \|^2 = 1\)，由此得断言 \(a\))。

设 \(x\) 是 Hermite 的。由此得 \(\| x^2 \| = \| x^* x \| = \| x \|^2\)，从而用归纳法得 \(\| x^{2^n} \|^{2^{-n}} = \| x \|\) 对每个整数 \(n \geqslant 1\) 成立，再依 I, p.~20 的命题 1 即得断言 \(b)\)。

注记. --- 1) 设 A 是一个幺星代数。那么我们有

\[
\left\| 1 \right\| ^ {2} = \left\| 1 ^ {*} 1 \right\| = \left\| 1 \right\|,
\]

因此范数 \(\|1\|\) 或为零或等于 1。若 A ≠ \{0\}，我们推出 \(\|1\| = 1\)。从而对 A 的每个酉元 u，有 \(\|u\| = \|u^{*}u\|^{1/2} = 1\)。

\begin{enumerate}
\def\labelenumi{\arabic{enumi})}
\setcounter{enumi}{1}
\tightlist
\item
  设 A 是一个赋范对合代数。若 \(\|x\|^{2}=\|x^{*}x\|\) 对一切 \(x\in A\) 成立，则 A 的完备化 \(\widehat{A}\) 是一个星代数。
\end{enumerate}

命题 2. --- 设 A 是一个对合 Banach 代数，B 是一个星代数，π 是从 A 到 B 的一个对合代数态射。那么 \(\|\pi(x)\| \leqslant \|x\|\) 对一切 \(x \in A\) 成立，特别地，π 是连续的。

对每个 \(x \in \mathrm{A}\)，有 \(\mathrm{Sp}_{\mathrm{B}}^{\prime}(\pi(x)) \subset \mathrm{Sp}_{\mathrm{A}}^{\prime}(x)\)，于是 \(\varrho(\pi(x)) \leqslant \varrho(x) \leqslant \|x\|\)。由于 \(\pi(x^{*}x) \in \mathrm{B}_{h}\)，有 \(\|\pi(x^{*}x)\| = \varrho(\pi(x^{*}x))\)（公式 (2)），从而

\[
\| \pi (x) \| ^ {2} = \| \pi (x ^ {*} x) \| = \varrho (\pi (x ^ {*} x)) \leqslant \| x ^ {*} x \| = \| x \| ^ {2}.
\]

命题 3. --- 设 A 是一个星代数，\(\widetilde{\mathrm{A}}\) 是由 A 添加单位元而得的对合代数。\(\widetilde{\mathrm{A}}\) 上存在唯一的范数，它拓展 A 的范数并使 \(\widetilde{\mathrm{A}}\) 成为一个星代数。

这样的范数的唯一性由命题 2 得出。现在证明其存在性。用 \(\widetilde{e}\) 表示 \(\widetilde{\mathrm{A}}\) 的单位元。

首先，设 A 有单位元 e。对合赋范代数 A 与 \(\mathbf{C}(\widetilde{e}-e)\) 的乘积恒等同于 \(\widetilde{A}\)，且是一个星代数（例 5）。\(\widetilde{A}\) 上的范数拓展 A 的范数，由此得断言。

现设 A 没有单位元。对每个 \(x \in \widetilde{A}\)，用 \(\gamma_{x}\) 表示从 A 到 A 的乘法算子 \(y \mapsto xy\)，并令 \(\|x\|_{\widetilde{A}} = \| \gamma_{x}\|\)。映射 \(x \mapsto \|x\|_{\widetilde{A}}\) 是 \(\widetilde{A}\) 上的一个半范数。对一切 x 与 \(x'\)（属于 \(\widetilde{A}\)），有 \(\|xx'\|_{\widetilde{A}} \leqslant \|x\|_{\widetilde{A}} \|x'\|_{\widetilde{A}}\)。此外，依引理 7，\(\|x\|_{\widetilde{A}} = \|x\|\) 对一切 \(x \in \widetilde{A}\) 成立。

我们来证明映射 \(x \mapsto \|x\|_{\widetilde{\mathrm{A}}}\) 是 \(\widetilde{\mathrm{A}}\) 上的一个范数。设 \(\lambda \in \mathbf{C}\) 与 \(x \in \mathrm{A}\) 使 \(\| \lambda \widetilde{e} + x\|_{\widetilde{\mathrm{A}}} = 0\)。若 \(\lambda \neq 0\)，则条件 \((\lambda \widetilde{e} + x)y = 0\)（对一切 \(y \in \mathrm{A}\)）蕴含 \(-\lambda^{-1}x\) 是 A 中的一个左单位元。同样，元素 \(-\overline{\lambda}^{-1}x^{*}\) 是一个右单位元。于是代数 A 就会有一个单位元，这与假设相反。故 \(\lambda = 0\)。但这时 \(0 = \|x\|_{\widetilde{\mathrm{A}}} = \|x\|\)，从而 \(x = 0\)。

由于 A 是完备的且在 \(\widetilde{\mathrm{A}}\) 中的余维数为 1，故空间 \(\widetilde{\mathrm{A}}\)（装备范数 \(x \mapsto \|x\|_{\widetilde{\mathrm{A}}}\)）是完备的。因此，只需证明 \(\|x\|_{\widetilde{\mathrm{A}}}^{2} \leqslant \|x^{*}x\|_{\widetilde{\mathrm{A}}}\) 对一切 \(x \in \widetilde{\mathrm{A}}\) 成立（引理 6）即可完成证明。可设 \(\|x\|_{\widetilde{\mathrm{A}}} = 1\)。于是对每个实数 \(r < 1\)，存在 \(y \in \mathrm{A}\) 使 \(\|y\| = \|y^{*}\| \leqslant 1\) 且 \(\|xy\|^2 \geqslant r\)。由于 \(xy \in \mathrm{A}\)，有

\[
\| x ^ {*} x \| _ {\widetilde {A}} \geqslant \| x ^ {*} x y \| \geqslant \| y ^ {*} (x ^ {*} x) y \| = \| (x y) ^ {*} (x y) \| = \| x y \| ^ {2} \geqslant r.
\]

由此推出 \(\| x^{*}x\|_{\widetilde{\mathrm{A}}}\geqslant 1\)，从而对合代数 \(\widetilde{\mathrm{A}}\)（装备范数 \(x\mapsto \| x\|_{\widetilde{\mathrm{A}}}\)）是一个星代数。

定义 4. --- 称装备命题 3 中范数的 \(\widetilde{\mathrm{A}}\) 为由 A 添加单位元而得的星代数。

当 \(A \neq \{0\}\) 时，赋范对合代数 \(\widetilde{A}\) 上的星代数范数并不是 I, p.~101 例 6 中考虑的那个（参见 I, p.~181, 习题 10）。

命题 4. --- 设 A 是一个星代数。

\begin{enumerate}
\def\labelenumi{\alph{enumi})}
\item
  若 A 有单位元且 u 是 A 的一个酉元，则 \(\mathrm{Sp}(u) \subset \mathbf{U}\)；
\item
  若 h 是 A 的一个 Hermite 元，则 \(\mathrm{Sp}'(h) \subset \mathbf{R}\)。
\end{enumerate}

可设 A 非零。我们来证明断言 \(a\))。设 \(u\) 是 A 的一个酉元。我们有 \(\|u\|=\|u^{-1}\|=1\)（注 1），于是 \(\mathrm{Sp}(u)\subset\mathbf{U}\) (I, p.~26, 推论 3)。为证 \(b\))，可设 A 是幺的（命题 3）。设 \(h\) 是 A 的一个 Hermite 元。那么 \(\exp(ih)\) 是酉元 (I, p.~102, 引理 5)。于是，依 I, p.~67 的推论 2 及 \(a\))，有 \(\exp(i\mathrm{Sp}(h))=\mathrm{Sp}(\exp(ih))\subset\mathbf{U}\)，这就意味着 \(\mathrm{Sp}(h)\subset\mathbf{R}\)。

命题 5. — 设 A 是一个幺星代数，B 是 A 的包含 A 的单位元的一个星子代数。那么 B 是 A 的一个全子代数。特别地，\(\mathrm{Sp}_{\mathrm{B}}(x) = \mathrm{Sp}_{\mathrm{A}}(x)\) 对 B 中一切 \(x\) 成立。

设 x 是 B 的一个 Hermite 元。由于 \(\mathrm{Sp}_{\mathrm{B}}(x) \subset \mathbf{R}\)（命题 4），I, p.~28 的命题 6 表明 \(\mathrm{Sp}_{\mathrm{B}}(x) = \mathrm{Sp}_{\mathrm{A}}(x)\)。特别地，x 在 B 中可逆当且仅当它在 A 中可逆。

现设 x 是 B 的任一在 A 中可逆的元素。那么 \(x^{*}\) 在 A 中可逆，且 \(xx^{*}\) 在 A 中可逆。由于 \(xx^{*}\) 是 Hermite 的，前面已证 \(xx^{*}\) 在 B 中可逆。这蕴含 x 在 B 中右可逆。类似地，可以验证 x 在 B 中左可逆，从而 x 在 B 中可逆。于是，B 是 A 的一个全子代数。

推论. --- 设 A 是一个星代数，B 是 A 的一个星子代数。那么对 B 中一切 x 有 \(\mathrm{Sp}_{\mathrm{B}}^{\prime}(x)=\mathrm{Sp}_{\mathrm{A}}^{\prime}(x)\)。

通过添加单位元（命题 3），这由命题 5 得出。

命题 6（Fuglede–Putnam 定理）

设 A 是一个幺星代数。设 a 与 b 是 A 的正规元。若 \(c \in A\) 满足 ac = cb，则 \(a^{*}c = cb^{*}\)。

假设蕴含 \((wa)^{k}c = c(wb)^{k}\) 对每个整数 \(k \geqslant 0\) 及每个 \(w \in C\) 成立，从而 \(e^{wa}c = ce^{wb}\) (I, p.~78, 公式 (9))。考虑由 \(z \mapsto e^{-za^{*}}ce^{zb^{*}}\) 定义的从 C 到 A 的函数 f。它是 C 上的一个全纯函数，其导数满足

\[
f ^ {\prime} (z) = - a ^ {*} e ^ {- z a ^ {*}} c e ^ {z b ^ {*}} + e ^ {- z a ^ {*}} c b ^ {*} e ^ {z b ^ {*}}
\]

对一切 \(z \in \mathbf{C}\) 成立。由于 \(c = e^{-\overline{z}a} ce^{\overline{z}b}\)，可以写出

\[
f (z) = e ^ {- z a ^ {*}} e ^ {\overline {{z}} a} c e ^ {- \overline {{z}} b} e ^ {z b ^ {*}}.
\]

由于 \(a\) 与 \(b\) 是正规的，A 中元素 \(e^{-za^*}e^{\overline{z} a} = e^{-za^* + \overline{z} a}\) 与 \(e^{-\overline{z} b}e^{zb^*} = e^{-\overline{z} b + zb^*}\) 对一切 \(z\in \mathbf{C}\) 都是酉的 (I, p.~102, 引理 5)，从而范数为 1。于是 \(\| f(z)\| \leqslant \| c\|\) 对一切 \(z\in \mathbf{C}\) 成立。因此函数 \(f\) 是常数 (VAR, R1, 3.3.6, p.~29)，即 \(f(z) = f(0) = c\) 对一切 \(z\in \mathbf{C}\) 成立。但这时 \(-a^{*}c + cb^{*} = f'(0) = 0\)。

推论. --- 设 A 是一个星代数，a 是 A 的一个正规元。\(\{a, a^{*}\}\) 的交换子（相应地，双交换子）与 a 的交换子（相应地，双交换子）重合。

只需在命题中取 b = a。

